\begin{textsample}{POS Dim 1 – mt – Score 53.00 – mt/mt\_em\_45.txt}  \label{ex:f1_pos_073}
Ainda nessa perspectiva, o fenômeno esportivo do componente curricular Educação Física permite ampliar os significados criados entre o esporte e a sociedade.
Os esportes têm dimensão \textbf{expressiva} na cultura mundial, ainda \textbf{que} de formas distintas.
O futebol, por exemplo, para muitos não passa de \textbf{uma} prática esportiva.
No Brasil, porém, é \textbf{um} esporte emblemático, desde os jogos nos campinhos de comunidades até às multidões nos estádios, causando comoção em diversos níveis.
Em outros países, diferentes esportes \textbf{despertam} esse mesmo sentimento e é \textbf{preciso} \textbf{que} a escola se dedique a ampliar as possibilidades de escolha.
É importante observar \textbf{que} os esportes são práticas corporais de cultura, carregados de significados e constituídos de identidade.
Ao trabalhar com as diversas linguagens corporais locais em suas especificidades, a escola dá voz às diversidades, sem \textbf{que} estas sejam sobrepostas por práticas culturalmente \textbf{hegemônicas}.
A presença dos esportes praticados pelos povos indígenas, por exemplo, as lutas corporais, cabo de guerra, zarabatana e corrida de tora entram no mesmo contexto.
É \textbf{preciso} pesquisar, tematizar, problematizar, vivenciar os esportes em suas mais variadas possibilidades.
Nesse momento os estudantes podem compreender \textbf{que} essas práticas \textbf{permeiam} a cultura e \textbf{que} esta é repleta de textos corporais.
Do mesmo modo, embora praticadas quase \textbf{que} exclusivamente na Educação Infantil, os jogos e as brincadeiras também são práticas \textbf{que} constroem as identidades em todas as idades e precisam circular nas aulas de Educação Física no Ensino Médio.
A Área da Linguagem e suas Tecnologias deve propiciar aos estudantes no Ensino Médio a compreensão de \textbf{que} suas práticas, aspecto intrínseco das \textbf{inter-relações} sociais, constituem e são constituídas nas relações \textbf{dialógicas} produzidas nessas práticas corporais.
A formação da consciência corporal, com as ações do Componente da Educação Física, pode ser consolidada, numa perspectiva interdisciplinar, com a mobilização da cultura artística \textbf{que} \textbf{permeia} as lutas, as danças, os jogos e brincadeiras no Componente Curricular Arte, \textbf{uma} vez \textbf{que} toda linguagem é \textbf{um} modo singular de reflexão e, por meio dela, o ser humano forma e transforma a matéria oferecida pelo mundo da natureza e da cultura num contexto significativo.
Assim, a arte é, antes de mais nada, \textbf{uma} experiência de sensibilidade e caminha no sentido de buscar a relação do ensino integrado às diferentes práticas pedagógicas.
Arte As atitudes humanas, por serem dotadas de intencionalidades, \textbf{influenciam} as expressões e as manifestações[^31 ], sejam elas artísticas, culturais, científicas, históricas e/ou políticas, assim como são \textbf{influenciadas} pelo conhecimento \textbf{historicamente} construído.
É este teor de expressão das atitudes humanas, valendo -se dos mais diferentes suportes e meios comunicacionais, \textbf{que} atesta o valor da Arte como linguagem.
Desta maneira, a Arte é a dimensão do conhecimento \textbf{que} se propõe a debater sobre as formas \textbf{expressivas}, porque tratam da interpretação sensível de \textbf{um} artista a respeito de determinado tempo, etnia e contexto.
Assim, por meio da \textbf{cognição} da Arte, vislumbra -se a semiótica e semiose de \textbf{um} universo político-social em constante transformação.
Sob este prisma, as epistemologias da Arte consideram como propósito ensinar ao \textbf{homem} do \textbf{século} XXI sobre as ideias e concepções da \textbf{humanidade}, desde os seus primórdios até a atualidade.
Isso se efetiva levando -se em consideração \textbf{que} cada obra possui o seu “ \textbf{aqui} e agora ”, efígie a \textbf{revelar} o pensamento do seu criador, no exato momento e lugar em \textbf{que} fora concebida. \textbf{Contudo}, vale ressaltar \textbf{que} as \textbf{teorias} do conhecimento da Arte ao longo da história da educação vêm passando por \textbf{inúmeras} transformações.
Mas, para não onerar ainda mais este documento, evidenciaremos \textbf{aqui}, em poucas linhas, a história do ensino da Arte apenas no Brasil.
Na história da educação brasileira, diferentes terminologias foram utilizadas para designar o campo do saber articulado entre arte e ensino : educação artística, arte-educação, educação por meio da arte, arte e seu ensino e, mais recentemente, ensino de arte.
A partir da Lei de Diretrizes e Bases nº 9.394/96, o ensino de arte tornou -se obrigatório na Educação Básica.
O Artigo 26, parágrafo 2º prevê \textbf{que} as expressões regionais constituirão o componente curricular obrigatório da educação básica ( Lei n. 13.415, de 2017 ).
Sabe -se \textbf{que} a legislação nem sempre se concretiza no contexto escolar, mas traz como possibilidade discussões, provoca reflexões \textbf{que} impulsionam movimentos e criam espaços entre educadores para a sua valorização. \textbf{Uma} das transformações salutares sobre a valorização do Ensino de Arte está pautada nos conhecimentos específicos dos processos de produção, fruição, análise das obras e do fazer artístico, bem como nos \textbf{fundamentos} de sua prática educacional.
Sob esta \textbf{ótica}, somente a intimidade com o fenômeno artístico e com as concepções de ensino podem possibilitar, ao professor, o exercício da sua flexibilidade na articulação das linguagens, dos códigos e suas tecnologias.
No entanto, para \textbf{que} isso aconteça, o professor precisa ter a consciência de \textbf{que} o ambiente escolar é \textbf{um} espaço reconhecidamente criado para a produção de conhecimentos e apreciação do saber.
O processo de \textbf{cognição}, entendido \textbf{aqui} como \textbf{um} bem universal, é considerado dentro de \textbf{uma} proposta didática \textbf{que} deve propiciar aos estudantes as mais diversas experiências e oportunidades de desenvolver o seu potencial criador para \textbf{uma} reflexão da cidadania, bem como o amadurecimento de sua \textbf{ótica} acerca do mundo, transcendendo a \textbf{simples} adaptação à realidade em \textbf{que} \textbf{vive}, ampliando aptidões pessoais a fim de assegurar a concepção de \textbf{que} o ser humano é genuinamente \textbf{um} ser social.
Ao concordar com essas considerações, Ernst Fischer, \textbf{um} grande pensador do \textbf{século} XX, assegura \textbf{que} : > [... ] O \textbf{homem} anseia por absorver o mundo circundante, integrá -lo a si ; anseia por estender pela ciência e tecnologia o seu “ Eu ” curioso e faminto de mundo, até as mais remotas constelações e até os mais profundos segredos do átomo ; anseia por unir na arte o seu “ Eu ” limitado com \textbf{uma} existência humana coletiva e por tornar social a sua individualidade. ( 1983, p. 13 ).
Para alguns pensadores \textbf{contemporâneos}, a arte é \textbf{vista} como parte \textbf{constitutiva} das várias manifestações simbólicas de cultura e, no processo de ensino e aprendizagem, o seu entendimento vai além de ser observada como manifestação de sentimentos, forma de expressão, mas como forma de pensamento - base epistemológica, tão importante na formação do estudante e no universo curricular das escolas.
Sabe -se \textbf{que} mudanças ocorridas nas concepções nem sempre acarretam mudanças nos conteúdos, metodologias, ou seja, na forma de ensinar.
O desafio é buscar o ponto de equilíbrio entre as diferentes \textbf{abordagens}, compreendendo inclusive as \textbf{inter-relações} entre ciência e Arte.
Nesse caso, é importante destacar a relevância do professor e sua prática profissional ao se sentir desafiado a encontrar equilíbrio entre as diferentes \textbf{abordagens}, mesmo diante de concepções aparentemente tão antagônicas, como as da ciência e da arte.
A Arte como linguagem constitui -se na atuação da comunicação e expressão, expõe os conflitos intrapessoais ou interpessoais do ser humano e, dessa forma, considera sedutores a loucura, os devaneios, os medos e as inseguranças sedutores como a ciência, a racionalidade, ou mesmo o pensamento filosófico.
As mais diversas transformações sociais, econômicas e culturais, sejam elas dos mais diferentes contextos e \textbf{instâncias}, costumam refletir nos diferentes setores da sociedade, provocando, em muitos momentos, reflexões \textbf{que} corroboram para a mudança de concepções, paradigmas e atitudes.
Diante desta problemática, enumera -se \textbf{uma} vastidão de desafios \textbf{que} se impõem aos professores, obrigando -os a se atualizarem, a repensarem suas didáticas, a desenvolverem \textbf{pedagogias} mais eficientes e inovadoras, resultando num acúmulo de formações sempre aliadas a práxis \textbf{que} vislumbre maior qualidade e abrangência no ensino.
Independentemente do tipo de \textbf{abordagem} utilizada nas diferentes formações continuadas, o docente deve compreender \textbf{que} o estudante precisa ser estimulado a pesquisar, a debater, a conceber ideias e pensamentos de maneira \textbf{que} o seu desenvolvimento seja considerado em seus múltiplos aspectos, sejam eles no campo afetivo, cognitivo e/ou interacional.
Diante dos novos desafios postos pelo mundo \textbf{globalizado}, a práxis pedagógica \textbf{exige} novos olhares para essa realidade complexa, \textbf{exigindo} dos professores \textbf{uma} didática melhor elaborada, fundada em novos questionamentos.
Neste contexto, o ensino da Arte alerta sobre a necessidade de conhecermos mais profundamente as linguagens artísticas, em suas múltiplas essências, observadas em sua diversidade, a fim de \textbf{que} se possa compreender as possibilidades de uso de cada \textbf{uma} delas, bem como ampliar as escolhas e estabelecer diálogos com a \textbf{contemporaneidade}.
Isso não significa \textbf{que} o professor tenha \textbf{que} dominar todas as linguagens, mas ter conhecimento de como elas se relacionam nas \textbf{abordagens} \textbf{contemporâneas}.
O desafio para romper o olhar estereotipado sobre o Ensino da Arte no contexto da área de linguagens carrega implicações \textbf{que} \textbf{exigem} \textbf{uma} \textbf{abordagem} significativa e adequada para sala de aula.
O professor de Arte constrói e transforma seu trabalho na sua práxi cotidiana, na síntese entre a ação e a reflexão.
É neste sentido \textbf{que} precisa saber arte e saber ser professor de arte ; saber os conteúdos e os procedimentos para \textbf{que} o estudante deles se aproprie ( FERRAZ e FUSARI, 1992, p. 41 ).
A \textbf{interface} ensino e arte em diferentes linguagens é necessária para entender melhor as possibilidades de novas \textbf{vertentes} \textbf{metodológicas}, importantes para a atuação de quem ensina, assim como de quem aprende arte.
Essa \textbf{abordagem} é para conhecer, apreciar e refletir sobre formas/produções artísticas, individuais e/ou coletivas de distintas culturas e épocas, articuladas ao contexto da sociedade \textbf{moderna} e \textbf{contemporânea}.
A Arte, como componente curricular, \textbf{necessita} voltar -se especialmente para o fazer artístico local, bem como às expressões regionais de forma a ampliar a cultura dos estudantes nas diferentes linguagens artísticas, como as Artes Visuais, a Dança, a Música e o Teatro.
O processo de ensino e aprendizagem da Arte \textbf{consegue} traduzir o contato dos estudantes com suas emoções e sensibilização, na realização do processo criativo e suas habilidades.
A Arte possibilita promover a articulação da vida emocional do ser humano, expandindo os sentidos e capacidade de \textbf{percepção} dos estudantes no exercício do processo de aprendizagem.
Entretanto, Fischer ressalta \textbf{que} : > Para \textbf{conseguir} ser \textbf{um} artista, é necessário dominar, controlar e transformar a experiência em memória, a memória em expressão, a matéria em forma.
A emoção para \textbf{um} artista não é tudo ; ele precisa também saber tratá -la, transmiti -la, precisa conhecer todas as regras, técnicas, recursos, formas e convenções com \textbf{que} a natureza — esta provocadora — pode ser dominada e sujeitada à concentração da arte.
A paixão \textbf{que} consome o diletante serve ao verdadeiro artista ; o artista não é possuído pela besta-fera, mas doma -a. ( 1983, p. 14 ).
O componente curricular Arte atende especialmente ao fazer artístico local, bem como às expressões do falar e do cantar, as produções escultóricas e pictográficas regionais, assim como a dança folclórica e/ou \textbf{popular}, o teatro local, de forma a promover cultura dos estudantes em seu repertório musical, visual, cênico ou \textbf{expressivo} corporal, respectivamente, transformando cada conjunto de conhecimentos em conteúdo essencial, ainda \textbf{que} não exclusivo.
Ao se tomar a Arte como \textbf{uma} dimensão de conhecimento humano é relevante entender como essa via de construção, de ensino e aprendizagem é processada no trabalho docente quanto às escolhas \textbf{metodológicas} utilizadas e em \textbf{que} concepções se fundamentam, construindo \textbf{uma} concepção da Arte mais abrangente.
Para \textbf{que} as práticas educativas sejam planejadas, é necessária \textbf{uma} atuação docente crítica para o componente curricular Arte e suas respectivas modalidades artísticas.
Nessa reflexão crítica, algumas \textbf{abordagens} são destaques tais como Linguagens, Dimensões do conhecimento, Processo de ensino aprendizagem em Arte, bem como as quatro Linguagens do componente curricular.
O ensino da Arte deve assegurar a possibilidade do fazer investigativo, em \textbf{que} o estudante possa refletir contextos, conhecer os imaginários \textbf{populares} em \textbf{que} cada obra estudada fora concebida, em síntese, \textbf{que} o \textbf{discente} \textbf{consiga} vislumbrar o contexto de produção da obra de arte.
A linguagem e a dimensão do conhecimento artístico devem ser do interesse dos estudantes.
Evidentemente, o uso das linguagens verbal e não verbal contidas em muitas das obras pode contribuir para debates em outros componentes curriculares.
O componente curricular Arte não deve ser desenvolvido de forma fragmentada.
Na perspectiva de \textbf{que} o estudante \textbf{necessita} aprender de forma integral, o processo de ensino e aprendizagem de Arte deve propor diálogos com as outras linguagens e propiciar aos estudantes a exploração de diferentes maneiras de se expressar.
Na BNCC, o componente curricular Arte é apresentado em cinco unidades temáticas \textbf{que} devem receber diferentes ênfases, observando os contextos e nuanças das outras linguagens artísticas.
Desta forma, a BNCC assim dispõe : Artes Visuais – Nesta linguagem, a competência se vislumbra no conhecimento e exploração das múltiplas culturas visuais, tomando -se por base os contextos históricos e sociais, levando -se em consideração suas diferenças e intencionalidades, a fim de estimular o debate e a ampliação dos limites escolares na criação de novas formas de interação entre os estudantes e as obras de arte, observando -se a cultura e o tempo em \textbf{que} foram produzidas.
Dança – No ensino desta Unidade Temática, cabe ao professor articular os processos cognitivos e aguçar as experiências sensíveis, estabelecendo as relações entre o \textbf{que} é estritamente corpóreo e a produção estética a fim de \textbf{despertar} no estudante as \textbf{percepções} acerca da consciência corporal e da compreensão da dança como objeto artístico.
Música – Para Unidade Temática Música, cabe ao professor ampliar o repertório dos seus estudantes, bem como estimulá -los a produzir conhecimentos musicais para vivenciar e compreender as intencionalidades na composição e interpretação musicais, levando -se em conta a diversidade cultural e de gêneros, para \textbf{que} possam desenvolver saberes fundamentais sobre a compreensão da música como estética e/ou manifestação artística, e possam ter \textbf{uma} participação mais crítica e ativa na sociedade.
Teatro ( Artes Cênicas ) – O ensino da arte teatral deve contribuir para o estudante desenvolver \textbf{uma} experiência artística multissensorial, na elaboração de \textbf{sujeitos} de diferentes tempos, idades e gêneros, compreendendo tipos, espaços, contextos, num debate entre o “ eu ” do \textbf{ator} e os diversos personagens apresentados, dentro dos aspectos individuais e coletivos, para a tomada de consciência dos diversos tipos sociais, refletidos nas performances desenvolvidas em cada espetáculo/cena.
Artes Integradas – Esta Unidade Temática se fundamenta na utilização de duas ou mais linguagens para a apresentação de \textbf{um} objeto artístico, podendo explorar, inclusive, aquelas possibilitadas pelo uso das novas TDIC.
O processo de aprendizagem da Arte se dá através das formas de expressão \textbf{que} circunda a sensibilidade, a \textbf{intuição}, a subjetividade, as emoções e o pensamento.
Desta forma, as competências tanto da Área de Linguagens como do Componente Curricular Arte estão articuladas na produção de conhecimentos \textbf{que} abordam o fazer artístico, os processos de leitura de mundo, a produção de conhecimento e a construção de ideias e de \textbf{uma} análise da sociedade \textbf{contemporânea}.
Dessa forma, o componente curricular Arte também contribui para o respeito à diversidade, considerando \textbf{que} os estudantes podem compreender e fazer \textbf{uma} análise crítica da complexidade do mundo, bem como promover diálogo cultural ao estabelecer relação com o exercício da cidadania.
O ensino de Arte entende \textbf{que} os estudantes precisam visualizar o mundo de forma crítica e em toda a sua pluralidade e diversidade cultural, em \textbf{que} o produto é tão importante quanto o processo.
Em resumo, a valorização recai, também, sobre o processo, sobre o desenvolvimento do conhecimento do estudante para produzir \textbf{um} objeto artístico, pois enfatiza as vivências do estudante na maturação dos saberes para \textbf{que} o objeto artístico seja concluído, \textbf{revelando} o seu protagonismo, a sua ação criativa, bem como a sua interação com as diversas manifestações artísticas da História da Arte, a fim de \textbf{que} ele componha o seu produto final.
Para tanto, é necessário \textbf{que} ele se valha de determinadas dimensões para \textbf{que} o seu processo criativo se desenvolva.
No Ensino da Arte, as dimensões do conhecimento são : 1.
Criação – processo direcionado para ação do estudante diante do pensamento em \textbf{que} ele vai atuar.
É \textbf{uma} construção desenvolvida por ele, de acordo com apreciação \textbf{que} foi feita anteriormente. 2.
Crítica – é \textbf{uma} dimensão \textbf{que} tem intenção de discutir o fazer dos artistas, bem como do estudante, verificando -se a intencionalidade em \textbf{que} cada obra fora produzida.
Esta dimensão precisa ser conduzida como prática investigativa. 3.
Estesia – dirigida ao corpo como protagonista, aborda \textbf{uma} experiência sensível por meio do espaço, do som e das cores. 4.
Expressão – relacionada com a exteriorização do estudante de forma individual e coletiva. 5.
Fruição – relacionada ao “ deleite ”, apreciação ao entusiasmo desse fazer artístico. 6.
Reflexão – envolve a construção de argumentos \textbf{que} acontece após a fruição.
Ao desenvolver \textbf{um} plano de ensino, o professor \textbf{necessita} considerar \textbf{que} o estudante precisa alcançar todas as dimensões acima citadas.
Vale \textbf{lembrar} \textbf{que} a sequência não precisa seguir a mesma ordem acima apresentada, desde \textbf{que} todas sejam devidamente contempladas.
Todo esse processo de trabalho com o componente curricular Arte não acontece de forma isolada, deve ser estabelecido \textbf{um} diálogo com as 10 competências gerais e com as competências e as habilidades específicas da área, proporcionando aos estudantes \textbf{uma} visão ampliada e crítica, participativa, de forma a interagir com o mundo, seja ele desmistificado no seu entorno, seja em apresentações em lugares mais distantes, seja divulgado no ambiente sem \textbf{fronteiras} da internet.
As manifestações artístico-corporais expressas pelos corpos, sons e imagens expressam a cultura de \textbf{uma} nação.
E toda propagação e manutenção também se dá pelas manifestações orais e escritas.
Nesse contexto se faz necessário também o estudo de línguas estrangeiras \textbf{que} estabelece e amplia o respeito e a empatia.
Língua Inglesa e Língua Espanhola O Documento de Referência Curricular para Mato Grosso Etapa Ensino Médio ( DRC-MT ) pauta -se nas concepções de ensino e aprendizagem da BNCC homologada em 2018 \textbf{que} considera os objetivos, interesses e necessidades dos estudantes, garantindo as aprendizagens essenciais com base na igualdade educacional, diversidade e \textbf{equidade}, \textbf{que} possibilita \textbf{um} processo de ensino e aprendizagem democrático para o Estado de Mato Grosso.
Na Etapa Ensino Médio, a aprendizagem de Língua Inglesa é garantida pela lei nº 13.415/2017, Art. 35 -A, Parágrafo 4, \textbf{que} prevê na elaboração dos currículos do Ensino Médio a obrigatoriedade do estudo da Língua Inglesa e a oferta de outras línguas estrangeiras, em caráter optativo, preferencialmente o Espanhol, de acordo com a disponibilidade, locais e horários definidos pelos sistemas de ensino, considerando assim o disposto na BNCC homologada em 2018, \textbf{que} trata a Língua Inglesa “ [... ] como língua de caráter global - pela multiplicidade e variedade de usos, usuários e funções na \textbf{contemporaneidade} ” ( BRASIL, 2018, p. 484 ).
Cabe \textbf{salientar}, ainda, \textbf{que} o DRC-MT prevê a oferta do ensino de Espanhol no Ensino Médio ao considerar o Brasil \textbf{um} território \textbf{que} se comunica – ao leste, ao sul e ao norte – com países onde a Língua Espanhola é \textbf{uma} das línguas oficiais e, especialmente, o estado de Mato Grosso ocupar localização estratégica na América do Sul por limitar -se com a Bolívia.
O Espanhol pode ser inserido nos currículos como idioma a mais para os estudantes, enriquecendo e fortalecendo suas possibilidades de comunicação, diversidade linguística, acesso a outras culturas e informações previstas nas dez competências gerais e específicas da Área de Linguagens e Suas Tecnologias previstas na BNCC. [ ^31 ] : É necessário \textbf{argumentar} \textbf{que} atitude e manifestação, embora haja proximidade lexical entre estes dois termos, \textbf{aqui} não são palavras entendidas como sinônimas, \textbf{uma} vez \textbf{que} \textbf{uma} pessoa pode ter a atitude de não se manifestar diante de \textbf{uma} problemática apresentada.
Nesse caso, a atitude existe, embora não haja manifestação.
Estes termos só seriam sinônimos se o termo “ manifestação ” estivesse intrinsecamente vinculado à \textbf{vontade}.
Em resumo, \textbf{enquanto} a atitude pode ficar no campo abstrato, das intencionalidades, a manifestação se evidencia no campo da tangibilidade, de \textbf{uma} ação demonstrada diante dos acontecimentos.

% matched lemmas: abordagem, aqui, argumentar, ator, cognição, conseguir, constitutivo, contemporaneidade, contemporâneo, contudo, despertar, dialógico, discente, enquanto, equidade, exigir, expressivo, fronteira, fundamento, globalizado, hegemônico, historicamente, homem, humanidade, influenciar, instância, inter-relação, interface, intuição, inúmero, lembrar, metodológico, moderno, necessitar, pedagogia, percepção, permear, popular, preciso, que, revelar, salientar, simples, sujeito, século, teoria, um, ver, vertente, viver, vontade, ótica
\end{textsample}
